\documentclass[12pt]{article}
\usepackage[margin=1in]{geometry}
\usepackage{amsmath,amssymb,amsthm,booktabs,longtable}
\usepackage{newtxtext,newtxmath}
\usepackage[hidelinks]{hyperref}
\newtheorem{proposition}{Proposition}
\newtheorem{lemma}{Lemma}
\theoremstyle{remark}
\newtheorem{remark}{Remark}
\newcommand{\E}{\mathbb E}
\renewcommand{\Pr}{\mathbb P}
\newcommand{\R}{\mathbb R}
\newcommand{\ind}{\mathbf 1}
\newcommand{\Dt}{\Delta_T}
\title{Entry Returns to Its Floor at the Strongest Incumbent:\\
Proposition 2(iii) of the Laplace Manuscript in Five Parts}
\author{}
\date{23 October 2026}
\begin{document}
\maketitle

\noindent
This note proves part (iii) of Proposition~2 of the Laplace manuscript
(\texttt{overleaf/paper}, commit \texttt{01bd5d6}) to the five-part standard of the team brief.
It uses the manuscript's own assumptions and notation. The manuscript proves part (iii) in
three sentences (Appendix A.2, ``Part (iii)''). They point to Steps~4 and~5 of its proof of
parts (i) and (ii). That argument is sound and assumes nothing it proves. It falls short of the
standard in two places. It does not display the equilibrium as one explicit tuple. It prints no
margin for its one new hypothesis at the benchmark. Both are supplied here. Numbers come from
\texttt{scripts/LA/margins.py} (run log \texttt{margins\_run.log}); every quantity that
involves $e$ is an outward interval enclosure at 50 decimal digits.

\section*{Setting}

The objects are those of the manuscript's Sections 2 to 3. The challenger's class is
$j\in\{H,L\}$ with prior $\pi=\Pr(j=H)\in(0,1)$; its value $\theta\sim G_j$ has a finite mean.
The incumbent's value $R\sim F$ is continuous on $[0,\bar r]$. The investor observes $j$ and
submits a (possibly mixed) order on $[-1,1]$ at linear cost $k|q|$. Noise $Z$ has the Laplace
density $f(z)=e^{-|z|/b}/(2b)$ with $b>1$; flow is $X=q+Z$. Market makers set
$P(X)=\E[V_T\mid X]$. The challenger sees $P$ and its cost $C\in\{0,c_H\}$, $\Pr(C=0)=\rho$,
then decides whether to prepare; an indifferent challenger prepares. Prepared bidders bid
truthfully in a second-price auction with reserve $p$. The auction objects are
$t_0=p[1-F(p)]$, $t_j$, $g_j$, $\Dt=t_H-t_L$, $w_L=t_L-t_0$, and $B_F(\mu)=g_L+\mu(g_H-g_L)$
(manuscript equations (6) and (9)). The posterior bounds are
\[
m(\pi)=\frac{\pi}{\pi+(1-\pi)e^{2/b}},\qquad
M(\pi)=\frac{\pi}{\pi+(1-\pi)e^{-2/b}},\qquad
\kappa(\pi)=\min\{m(\pi),\,1-M(\pi)\}.
\]

\paragraph{Hypotheses.}
Fix the incumbent distribution $F_2$ (the manuscript's strongest incumbent) and suppose:
\begin{enumerate}
\item[(H1)] \emph{Classes.} $G_H(u)\le G_L(u)$ for every $u$, and $G_H(p^-)=G_L(p^-)=0$
  (the restrictions of the manuscript's Section 2.2).
\item[(H2)] \emph{Incumbent.} $F_2$ is continuous on $[0,\bar r_2]$ and extended by one above.
\item[(H3)] \emph{Spread.} $\Dt(F_2)>0$.
\item[(H4)] \emph{Trading cost.} $0<k<(1-1/b)\,\rho\,\kappa(\pi)\,\Dt(F_2)$.
\item[(H5)] \emph{Exclusion.} $B_{F_2}(M(\pi))<c_H$.
\item[(H6)] $b>1$, $0<\rho<1$, $0<\pi<1$, $c_H>0$.
\end{enumerate}
In the manuscript, (H3) and (H4) are not separate
assumptions. There $F_2$ is a first-order stochastic strengthening of $F_1$, so Proposition~1(b)
gives $\Dt(F_2)\ge\Dt(F_1)$, and (A2) at $F_1$ gives both. (H5) is the one new hypothesis of part (iii). Under (H1), equation (6) gives
$g_H\ge g_L\ge0$ and $\Dt\ge0$; (H3) sharpens the last to a strict inequality.

\section{The equilibrium object}

The candidate is the tuple $\mathcal E^*=(\sigma_H^*,\sigma_L^*;\,\mu^*;\,P^*;\,\beta^*;\,e^*;\,\text{bids})$:
\begin{enumerate}
\item \emph{Orders.} $\sigma_H^*=\delta_{+1}$ and $\sigma_L^*=\delta_{-1}$: the investor buys one
  unit in class $H$ and sells one unit in class $L$.
\item \emph{Market posterior.} For every $x\in\R$,
  \[
  \mu^*(x)=\frac{\pi f(x-1)}{\pi f(x-1)+(1-\pi)f(x+1)}
  =\Big[1+\frac{1-\pi}{\pi}\exp\Big\{-\frac{|x+1|-|x-1|}{b}\Big\}\Big]^{-1},
  \]
  equal to $m(\pi)$ on $x\le-1$, strictly increasing on $(-1,1)$ with log odds
  $\operatorname{logit}\pi+2x/b$, and equal to $M(\pi)$ on $x\ge1$.
\item \emph{Price function.} For every $x\in\R$,
  \[
  P^*(x)=t_0+\rho\,[\,w_L+\Dt\,\mu^*(x)\,].
  \]
  Its image is the closed interval $[P_-,P_+]$ with $P_\pm=t_0+\rho[w_L+\Dt M_\pm]$,
  $M_-=m(\pi)$, $M_+=M(\pi)$. The endpoints are atoms: $P^*=P_-$ on $\{X\le-1\}$ and
  $P^*=P_+$ on $\{X\ge1\}$. Their probabilities are $\Pr(X\le-1\mid H)=\tfrac12e^{-2/b}$,
  $\Pr(X\le-1\mid L)=\tfrac12$, $\Pr(X\ge1\mid H)=\tfrac12$, $\Pr(X\ge1\mid L)=\tfrac12e^{-2/b}$.
\item \emph{Challenger's belief.} At a realized price $P\in[P_-,P_+]$,
  \[
  \beta^*(P)=\frac{P-t_0-\rho w_L}{\rho\Dt}\in[m(\pi),M(\pi)].
  \]
  No other price is reached on path or after any unilateral investor deviation (Section~3.3), so
  no further belief is specified.
\item \emph{Entry rule.} At every price and cost, $e^*(P,0)=1$ and $e^*(P,c_H)=0$. The
  cost-averaged entry rate is $e^*(P)=\rho$ at every price.
\item \emph{Bids.} Each prepared bidder bids its value; a prepared challenger with
  $\theta<p$ does not bid.
\end{enumerate}

\section{The conditions}

$\mathcal E^*$ is an equilibrium if the following hold.
\begin{enumerate}
\item[(C1)] \emph{Bayes consistency.} $\mu^*(x)=\Pr(j=H\mid X=x)$ under $(\sigma_H^*,\sigma_L^*)$,
  and $\beta^*(P^*(X))=\Pr(j=H\mid P^*(X))$ almost surely.
\item[(C2)] \emph{Competitive prices.} $P^*(x)=\E[V_T\mid X=x]$ for every $x$, where $V_T$ is
  computed under $\beta^*$, $e^*$ and truthful bids.
\item[(C3)] \emph{Entry optimality.} At every reached price and each cost, the challenger's
  preparation decision maximizes expected gross profit net of cost, with the tie rule.
\item[(C4)] \emph{Investor optimality.} For each class $j$, the order $q_j^*$ maximizes
  $q\,\E[V_T-P^*(X)\mid j,\,X=q+Z]-k|q|$ over $q\in[-1,1]$, and no mixture does better,
  with $P^*$ and $e^*$ held fixed.
\item[(C5)] \emph{Truthful bids.} Bidding one's value is weakly dominant.
\end{enumerate}

\section{Verification against the fixed object}

\subsection{Bayes consistency (C1)}
Under $\sigma^*$ the conditional flow densities are $f(x-1)$ in class $H$ and $f(x+1)$ in
class $L$. Bayes' rule gives the first form of $\mu^*$. Since
$\log[f(x-1)/f(x+1)]=(|x+1|-|x-1|)/b$, which equals $-2/b$ for $x\le-1$, $2x/b$ for
$-1<x<1$, and $2/b$ for $x\ge1$, the three-piece description follows. For the price
posterior: $P^*$ is the composition of $\mu^*$ with the strictly increasing affine map
$\mu\mapsto t_0+\rho[w_L+\Dt\mu]$ (slope $\rho\Dt>0$ by (H3), (H6)). Hence
$P^*(x)=P^*(x')$ if and only if $\mu^*(x)=\mu^*(x')$, and $\mu^*(X)$ is a function of
$P^*(X)$, namely $\mu^*(X)=\beta^*(P^*(X))$. The tower property gives
$\Pr(H\mid P^*(X))=\E[\mu^*(X)\mid P^*(X)]=\beta^*(P^*(X))$. At the two atoms the preimage
is a half-line on which $\mu^*$ is constant, so conditioning on the whole preimage gives the
same value.

\subsection{Competitive prices (C2)}
Conditional on class $j$ and flow $x$, the price $P^*(x)$ is fixed, the cost $C$ is independent
of $(j,\theta,R,Z)$, and the entry rule is $e^*$. The challenger prepares exactly when $C=0$,
with probability $\rho$. Given preparation and class $j$, expected target proceeds are $t_j$: the
challenger's value is independent of flow given its class, and bids are truthful. Without
preparation they are $t_0$. Hence
\[
\E[V_T\mid j,\,X=x]=t_0+\rho\,(t_j-t_0),
\]
and averaging over the class with $\mu^*(x)$,
\[
\E[V_T\mid X=x]=t_0+\rho\big[\mu^*(x)(t_H-t_0)+(1-\mu^*(x))(t_L-t_0)\big]
=t_0+\rho\,[\,w_L+\Dt\mu^*(x)\,]=P^*(x).
\]

\subsection{Entry optimality (C3)}
At price $P\in[P_-,P_+]$ the challenger's belief is $\beta^*(P)\in[m(\pi),M(\pi)]$. Its
expected gross profit is $B_{F_2}(\beta^*(P))$, because its value depends on the price only
through its class. Two facts about $B_{F_2}$:
\begin{itemize}
\item $B_{F_2}(\mu)\ge g_L(F_2)\ge0$ for every $\mu\in[0,1]$, by (H1) and (6); the gross
  profit is an expectation of nonnegative auction rents.
\item $B_{F_2}$ is nondecreasing ($g_H\ge g_L$), so $B_{F_2}(\mu)\le B_{F_2}(M(\pi))<c_H$ for
  every $\mu\le M(\pi)$, by (H5).
\end{itemize}
The zero-cost challenger compares a nonnegative gross profit with a zero cost: it prepares,
strictly if $g_L(F_2)>0$ and by the tie rule otherwise. The high-cost challenger compares a
gross profit below $c_H$ with the cost $c_H$: it stays out, strictly. This is $e^*$.

\subsection{Investor optimality (C4)}
Fix $P^*$ and $e^*$. Consider class $H$ and an order $q\in[-1,1]$, so $X=q+Z$. Conditional on
$(H,X=x)$ the terminal value has mean $t_0+\rho(t_H-t_0)$ by Section~3.2. The entry rate is
$\rho$ at every price, so the investor's order changes nothing in the challenger's decision.
The residual advantage at flow $x$ is therefore
\[
A_H(x)=t_0+\rho(t_H-t_0)-P^*(x)=\rho\,\Dt\,[1-\mu^*(x)],\qquad
A_L(x)=P^*(x)-t_0-\rho(t_L-t_0)=\rho\,\Dt\,\mu^*(x),
\]
the second for class $L$. Both are bounded and measurable, and
\[
\rho\,\kappa(\pi)\,\Dt\;\le\;A_H(x),\,A_L(x)\;\le\;\Dt\qquad\text{for every }x,
\]
since $\mu^*\in[m(\pi),M(\pi)]$ gives $1-\mu^*\ge1-M(\pi)\ge\kappa(\pi)$ and
$\mu^*\ge m(\pi)\ge\kappa(\pi)$.

A correctly signed order of magnitude $s\in[0,1]$ ($q=s$ in class $H$, $q=-s$ in class $L$)
earns
\[
U_j(s)=s\,\bar A_j(s)-ks,\qquad
\bar A_H(s)=\int_\R f(x-s)A_H(x)\,dx,\qquad
\bar A_L(s)=\int_\R f(x+s)A_L(x)\,dx .
\]

\begin{lemma}[translation derivative]\label{lem:w11}
Let $A:\R\to[0,\bar A]$ be measurable, $\epsilon\in\{-1,+1\}$, and
$\bar A_\epsilon(s)=\int f(x-\epsilon s)A(x)\,dx$ for $s\in[0,1]$. Then $\bar A_\epsilon$ is
absolutely continuous on $[0,1]$, and for almost every $s$
\[
\bar A_\epsilon'(s)=-\epsilon\int f'(x-\epsilon s)A(x)\,dx,\qquad
|\bar A_\epsilon'(s)|\le\frac1b\,\bar A_\epsilon(s).
\]
\end{lemma}
\begin{proof}
The Laplace density is absolutely continuous on $\R$ with derivative
$f'(z)=-\operatorname{sgn}(z)f(z)/b$ for $z\ne0$, so $|f'|\le f/b$ almost everywhere and
$\int|f'|=1/b$. For $0\le s_1<s_2\le1$ and every $x$,
\[
f(x-\epsilon s_2)-f(x-\epsilon s_1)=-\epsilon\int_{s_1}^{s_2}f'(x-\epsilon u)\,du .
\]
Multiply by $A(x)$ and integrate over $x$. The double integral of the absolute value is at
most $\bar A\,(s_2-s_1)\int|f'|=\bar A(s_2-s_1)/b<\infty$. Fubini's theorem therefore allows
the order of integration to be exchanged:
\[
\bar A_\epsilon(s_2)-\bar A_\epsilon(s_1)=\int_{s_1}^{s_2}\Big[-\epsilon\int f'(x-\epsilon u)A(x)\,dx\Big]du .
\]
The bracket is an integrable function of $u$, so $\bar A_\epsilon$ is absolutely continuous
and its derivative equals the bracket almost everywhere. Finally, $A\ge0$ and $|f'|\le f/b$
give $|\int f'(x-\epsilon s)A(x)\,dx|\le b^{-1}\int f(x-\epsilon s)A(x)\,dx=\bar A_\epsilon(s)/b$.
\end{proof}

Apply the lemma with $A=A_H$, $\epsilon=+1$ and with $A=A_L$, $\epsilon=-1$. Then $U_j$ is
absolutely continuous on $[0,1]$ and, for almost every $s$,
\[
U_j'(s)=\bar A_j(s)+s\bar A_j'(s)-k\;\ge\;\Big(1-\frac sb\Big)\bar A_j(s)-k
\;\ge\;\Big(1-\frac1b\Big)\rho\,\kappa(\pi)\,\Dt-k\;=:\;\zeta_{\mathrm{tr}}>0,
\]
The second inequality uses $1-s/b\ge1-1/b>0$, by (H6), and
$\bar A_j(s)\ge\rho\kappa(\pi)\Dt$, because the pointwise lower bound integrates against a
probability density. The last inequality is (H4). Integrating,
\[
U_j(1)-U_j(s)=\int_s^1U_j'(u)\,du\;\ge\;\zeta_{\mathrm{tr}}\,(1-s)>0\qquad\text{for every }s\in[0,1).
\]
So the unit correctly signed order strictly dominates every smaller correctly signed
magnitude, including zero. A wrongly signed order of magnitude $s>0$ earns
$-s\int f(x+s)A_H(x)\,dx-ks<0=U_H(0)$ in class $H$, and symmetrically in class $L$, because
the residuals are nonnegative. Hence $q_H^*=1$ is the unique maximizer in class $H$ and
$q_L^*=-1$ in class $L$. A mixture that puts positive mass on any other order averages
strictly smaller payoffs and is strictly worse. This verifies (C4).

\subsection{Truthful bids (C5)}
In a second-price auction with reserve, a bid equal to one's value is weakly dominant against
every competing bid. A higher bid can only add unprofitable wins. A lower bid can only discard
profitable wins. The winner's payment does not depend on its own bid. The price and the entry
decision change beliefs about the competing bid, not this pointwise comparison.

\subsection{Outcome at $F_2$}
Total entry is $\E[e^*(P^*(X))]=\rho$. Class-$H$ ownership is $\pi\rho W_H(F_2)$ with
$W_H(F)=\E_H[F(\theta)\ind\{\theta\ge p\}]$. The price is informative: $\mu^*$ is nonconstant.

\section{Uniqueness through bounds on every profile}

Let $\mathcal E=(\sigma_H,\sigma_L;\,\mu;\,P;\,\beta;\,e;\,\text{bids})$ be any equilibrium
of the game at $F_2$, with arbitrary mixed orders. Each step below is a bound that holds for
this arbitrary profile; none uses the candidate $\mathcal E^*$.

\paragraph{Step 1: posterior bounds for every profile.}
Write $a_j(x)=\int f(x-q)\,\sigma_j(dq)$. For any $q,q'\in[-1,1]$ the triangle inequality gives
$\big||x-q|-|x-q'|\big|\le|q-q'|\le2$, so $e^{-2/b}\le f(x-q)/f(x-q')\le e^{2/b}$. Integrating
against $\sigma_H(dq)\sigma_L(dq')$ gives $e^{-2/b}a_L(x)\le a_H(x)\le e^{2/b}a_L(x)$ for every
$x$. By Bayes' rule the log odds of $\mu(x)=\pi a_H/(\pi a_H+(1-\pi)a_L)$ lie within $2/b$ of
$\operatorname{logit}\pi$, so $m(\pi)\le\mu(x)\le M(\pi)$ for every $x$. The price is a
function of $X$, so $\beta(P(X))=\E[\mu(X)\mid P(X)]$ lies in the same interval almost surely.
Both conditional flow densities are positive on $\R$. The flow law under every pure or mixed
order, and under every unilateral deviation, is therefore equivalent to Lebesgue measure.

\paragraph{Step 2: entry is $\rho$ for every profile.}
At every reached price the challenger's belief $\beta(P)$ lies in $[m(\pi),M(\pi)]$ by Step~1,
and its gross profit is $B_{F_2}(\beta(P))$. Section~3.3 applies verbatim: the zero-cost
challenger prepares (strictly or by the tie rule), and the high-cost challenger stays out
because $B_{F_2}(\beta(P))\le B_{F_2}(M(\pi))<c_H$. Hence the cost-averaged entry rate is
$e(P)=\rho$ at every reached price, in every equilibrium. This step does not need the price to
reveal the posterior.

\paragraph{Step 3: price and residuals for every profile.}
With $e\equiv\rho$ and $C$ independent of everything else, Section~3.2 gives
$\E[V_T\mid j,X=x]=t_0+\rho(t_j-t_0)$. Competitive pricing then gives
$P(x)=t_0+\rho[w_L+\Dt\mu(x)]$ for almost every $x$, where $\mu$ is this profile's own
posterior. The residuals under a unilateral deviation, which holds $P$ and $e$ fixed, are
$A_H(x)=\rho\Dt[1-\mu(x)]$ and $A_L(x)=\rho\Dt\mu(x)$. Step~1 gives
$\rho\kappa(\pi)\Dt\le A_H,A_L\le\Dt$ for every $x$. The pricing identity holds for almost
every flow. Its exceptional set is Lebesgue-null. Every order gives flow a density, so the set has
probability zero under every deviation and leaves the deviation payoff unchanged.

\paragraph{Step 4: full orders for every profile.}
Lemma~\ref{lem:w11} and the bounds of Step~3 give
$U_j'(s)\ge\zeta_{\mathrm{tr}}>0$ almost everywhere on $[0,1]$ for both classes, as in
Section~3.4. The bound holds against the fixed $P$ and $e$ of this profile. Therefore the unit correctly signed order is the unique
best response in each class, and $\sigma_H=\delta_{+1}$, $\sigma_L=\delta_{-1}$. Mixing on any
other order is strictly worse.

\paragraph{Step 5: the rest of the tuple.}
With the orders pinned, $\mu=\mu^*$ everywhere by Bayes' rule and $P=P^*$ for almost every
flow by Step~3. Also $\beta(P^*(X))=\mu^*(X)$ almost surely by (C1), $e\equiv\rho$ by
Step~2, and bids are truthful. Every equilibrium therefore has the trading outcome $(1,-1)$,
the on-path entry $\rho$, and the price $P^*$ up to a flow-null set. This is the manuscript's
uniqueness scope: trading and on-path entry under truthful bidding, with arbitrary mixed orders
and every unilateral deviation. Prices at flows in a null set, and beliefs at prices outside
$[P_-,P_+]$, are not pinned. Neither is reached on path or by any deviation: every order gives
flow a positive density on $\R$, and $P^*$ maps $\R$ onto $[P_-,P_+]$.

\paragraph{The return to the floor.}
Parts (i) and (ii) of the manuscript's Proposition~2 give, at $F_1$ under (A1) and (A2), the
unique trading outcome $(1,-1)$ with entry
$\mathsf E_1=\rho+(1-\rho)[\pi\alpha_H+(1-\pi)\alpha_L]>\rho$. Entry exceeds $\rho$ there
for two reasons. The threshold belief $\tau_1=(c_H-g_L(F_1))/(g_H(F_1)-g_L(F_1))$ lies strictly
below $M(\pi)$. The flow tail $\{X\ge1\}$, on which $\mu^*=M(\pi)$, has positive probability
in both classes. At $F_2$ the
trading outcome is again $(1,-1)$ and entry is exactly $\rho$. Entry rises from $\rho$ at
$F_0$ to $\mathsf E_1$ at $F_1$ and returns to $\rho$ at $F_2$. The price stays informative at
$F_2$. This is the stated rise and fall across three economies; it says nothing about the
strengths between them.

\begin{proposition}[Proposition 2(iii), restated; analytical]
Under (H1)--(H6), $\mathcal E^*$ is an equilibrium. Every equilibrium has trading outcome
$(q_H,q_L)=(1,-1)$, on-path entry $\rho$, and the price $P^*$ up to a flow-null set.
\end{proposition}

\section{Strict margins and the open set}

Every hypothesis is a strict inequality. Take the benchmark of manuscript Appendix A.8
($h=10$, $\ell=1$, $p=1/2$, $\rho=1/4$, $c_H=6$, $b=2$, $k=1/50$, $\pi=1/2$) with
$r_2=18/5$. The point-mass auction objects are then exact rationals:
$g_H(r_2)=5879/720$, $g_L(r_2)=5/48$, $\Dt(r_2)=169/180$, $\tau(r_2)=4245/5804$, and
$M=e/(1+e)$, $m=1/(1+e)$, $\kappa=m$. The margins are:

\begin{center}
\begin{tabular}{llll}
\toprule
Hypothesis & Margin & Value at B & Type \\
\midrule
(H5) exclusion & $c_H-B_{r_2}(M)$ & $[0.00268890226590,\ 0.00268890226591]$ & interval \\
(H5) in posterior units & $\tau(r_2)-M$ & $[0.000333564719409,\ 0.000333564719410]$ & interval \\
(H4) trading bound & $(1-1/b)\rho\kappa\Dt(r_2)-k$ & $[0.0115632640357,\ 0.0115632640358]$ & interval \\
(H3) spread & $\Dt(r_2)$ & $169/180=0.93888\ldots$ & exact \\
(H3) via Prop.\ 1(b) & $\Dt(r_2)-\Dt(r_1)$ & $49/180=0.27222\ldots$ & exact \\
(H6) score bound & $b-1$ & $1$ & exact \\
(H6) floor & $\rho$, $1-\rho$ & $1/4$, $3/4$ & exact \\
Support & $\ell-p$, $r_2-\ell$, $h-r_2$ & $1/2$, $13/5$, $32/5$ & exact \\
Zero-cost strictness & $g_L(r_2)$ & $5/48$ & exact \\
Strengthening & $r_2-r_1$ & $3/5$ & exact \\
Position above the ceiling & $r_2-r_C$ & $[0.00734148120533,\ 0.00734148120534]$ & interval \\
\bottomrule
\end{tabular}
\end{center}

\noindent
Here $r_C=[Mh-c_H+\sqrt{(Mh-c_H)^2+M((1-M)\ell^2-p^2)}]/M$ is the manuscript's ceiling
(A.11), enclosed as $[3.592658518794,\ 3.592658518795]$. The function $r\mapsto B_r(M)$ has
derivative $-M/2-[(1-M)\ell^2-p^2]/(2r^2)$. At B it equals $-0.36626\ldots$. It is negative on
all of $r>\ell>p$: if $(1-M)\ell^2\ge p^2$ both terms are negative, and otherwise
$Mr^2>M\ell^2>p^2-(1-M)\ell^2$ because $\ell>p$. Hence (H5) holds at the benchmark exactly
when $r_2>r_C$.

\paragraph{The open set.}
At the benchmark, fix all coordinates but one. The hypotheses of this note, together with (A1)
and (A2) at $(r_0,r_1)$, then hold on the following open sets:
\begin{itemize}
\item $r_2\in(r_C,h)=(3.59266\ldots,\,10)$; the benchmark $3.6$ lies $0.00734$ above the lower end.
\item $c_H\in(B_{3.6}(M),\,B_3(M))=(5.99731\ldots,\,6.21715\ldots)$; the benchmark $6$ lies
  $0.00269$ above the lower end. (The lower bound of (A1), $B_{1.2}(1/2)=4.804$, does not bind.)
\item $\pi\in(0.467850\ldots,\,0.500424\ldots)$; the lower end is where $c_H=B_3(M(\pi))$,
  the upper end is where $B_{3.6}(M(\pi))=c_H$. The benchmark $1/2$ lies $0.000424$ below the
  upper end.
\item $k\in(\Dt(1.2),\,(1-1/b)\rho m\Dt(3))=(0.016667,\,0.022412)$; (H4) at $r_2$ then holds
  with the larger margin above.
\end{itemize}
All inequalities are strict and all quantities are continuous in $(r_0,r_1,r_2,c_H,k,\pi,\rho,b)$,
so the joint region is open and contains the benchmark. The three thinnest margins are
(H5) in $r_2$, (H5) in $c_H$, and the prior window's upper end. Each is of order $10^{-3}$ or
smaller relative to the benchmark value. The result is robust to perturbation, but only inside
this narrow band at $r_2=3.6$. The manuscript prints $r_C$ and the prior endpoint $0.500424$
but frames neither as a margin; it never prints $c_H-B_{3.6}(M)$.

\paragraph{Status.}
The proposition is analytical. The margins table is computer-assisted. The rational rows use
exact rational arithmetic. The rows that involve $e$ use outward interval arithmetic at 50
digits. The script is \texttt{scripts/LA/margins.py}, with run log
\texttt{scripts/LA/margins\_run.log} and output \texttt{scripts/LA/margins.csv}.

\paragraph{Tie rules.}
The challenger's rule (an indifferent challenger prepares) is not used at the benchmark.
There $g_L(r_2)=5/48>0$ makes the zero-cost decision strict, and (H5) makes the high-cost
decision strict. The rule would be used only if $g_L(F_2)=0$. The investor's best responses are
strict, so no investor tie rule is needed.

\paragraph{Verdict on the manuscript's proof.}
The manuscript's argument for part (iii) is the one given here in compressed form. The
posterior bound and the floor hold for every profile. (H5) removes high-cost entry at every
reachable belief. Proposition~1(b) carries the trading bound from $F_1$ to $F_2$. The derivative
bound forces full orders. It is non-circular. It lacks three things against the five-part
standard. These are the explicit tuple (Section~1); the observation that uniqueness in part
(iii) needs no price inversion (Step~2); and the margins table (Section~5). The manuscript's
printed statement of uniqueness scope matches Step~5.

\end{document}
